\subsection{Expresiones Lambda}

\subsubsection{Regla: solo para operaciones concisas de una línea}

\textbf{Regla:} Las lambdas (\texttt{=>}) solo para operaciones que quepan en una línea. Prohibido un cuerpo de bloque con varias instrucciones (Martin, 2008).

\textbf{Código correcto:}
\begin{lstlisting}[style=csharpstyle]
_players.RemoveAll(player => !player.IsActive);
\end{lstlisting}

\textbf{Código incorrecto:}
\begin{lstlisting}[style=csharpstyle]
_players.RemoveAll(player =>
{
    bool shouldRemove = !player.IsActive;
    _logger.Debug($"Player evaluated. Player={player.Name}.");
    return shouldRemove;
});
\end{lstlisting}

\subsubsection{Regla: contextos en los que se prohíbe una expresión lambda}

\textbf{Regla:} Se prohíbe la expresión lambda en tres contextos:

\begin{itemize}
  \item \textbf{Evento cuya suscripción deba cancelarse.} Cada aparición de una lambda crea un delegado distinto, por lo que \texttt{-=} no retira lo que registró \texttt{+=}; el manejador se declara como método con nombre (Microsoft, 2015).
  \item \textbf{Asignación sobre el estado capturado.} El cuerpo de una lambda solo lee sus entradas y devuelve un resultado; para modificar una variable capturada o un campo se usa \texttt{foreach} o un método con nombre (Microsoft, 2026g).
  \item \textbf{Lambda dentro de otra lambda.} La lambda interior se extrae a un método privado y se pasa como grupo de método (Sección 5.2).
\end{itemize}

\textbf{Descripción:} Las lambdas se retiran del código sobre todo por pérdida de legibilidad, dificultad para depurar y degradación del rendimiento (Zheng et al., 2021). Un manejador registrado con lambda reúne las tres y, además, mantiene vivo al suscriptor mientras exista el publicador.

\textbf{Código correcto:}
\textit{Suscripción a un evento:}
\begin{lstlisting}[style=csharpstyle]
public sealed class ScoreBoardView
{
    private readonly ITurnManager _turnManager;

    public ScoreBoardView(ITurnManager turnManager)
    {
        _turnManager = turnManager;
    }

    public void Attach()
    {
        _turnManager.TurnChanged += HandleTurnChanged;
    }

    public void Detach()
    {
        _turnManager.TurnChanged -= HandleTurnChanged;
    }

    private void HandleTurnChanged(object? sender, TurnChangedEventArgs e)
    {
        Render(e.CurrentPlayer);
    }

    private void Render(Player currentPlayer)
    {
    }
}
\end{lstlisting}

\textit{Asignación sobre el estado capturado:}
\begin{lstlisting}[style=csharpstyle]
public int CalculateTotalScore()
{
    return _players.Sum(player => player.Score);
}
\end{lstlisting}

\textit{Lambda dentro de otra lambda:}
\begin{lstlisting}[style=csharpstyle]
public IReadOnlyList<Player> FindPlayersWithWinningCard()
{
    return _players.Where(HasWinningCard).ToList();
}

private bool HasWinningCard(Player player)
{
    return player.Cards.Any(card => card.IsWinning);
}
\end{lstlisting}

\textbf{Código incorrecto:}
\textit{Suscripción a un evento:}
\begin{lstlisting}[style=csharpstyle]
public sealed class ScoreBoardView
{
    private readonly ITurnManager _turnManager;

    public ScoreBoardView(ITurnManager turnManager)
    {
        _turnManager = turnManager;
    }

    public void Attach()
    {
        _turnManager.TurnChanged += (sender, e) => Render(e.CurrentPlayer);
    }

    public void Detach()
    {
        _turnManager.TurnChanged -= (sender, e) => Render(e.CurrentPlayer);
    }

    private void Render(Player currentPlayer)
    {
    }
}
\end{lstlisting}

\textit{Asignación sobre el estado capturado:}
\begin{lstlisting}[style=csharpstyle]
public int CalculateTotalScore()
{
    int totalScore = 0;
    _players.ForEach(player => totalScore += player.Score);
    return totalScore;
}
\end{lstlisting}

\textit{Lambda dentro de otra lambda:}
\begin{lstlisting}[style=csharpstyle]
public IReadOnlyList<Player> FindPlayersWithWinningCard()
{
    return _players.Where(player => player.Cards.Any(card => card.IsWinning)).ToList();
}
\end{lstlisting}
